\documentclass[12pt]{article}
\usepackage[margin=1in]{geometry}
\usepackage{amsmath,amssymb,amsthm}
\usepackage{booktabs}
\usepackage[colorlinks=true,linkcolor=blue,citecolor=blue,urlcolor=blue]{hyperref}

\newcommand{\fl}[1]{\lfloor #1 \rfloor}
\newcommand{\ce}[1]{\lceil #1 \rceil}
\newcommand{\g}{\gamma_{\{2\}}}
\newcommand{\G}[2]{P_{#1}\,\square\,P_{#2}}
\newcommand{\N}{\mathbb{N}}
\newcommand{\Z}{\mathbb{Z}}

\theoremstyle{plain}
\newtheorem{theorem}{Theorem}
\newtheorem{lemma}[theorem]{Lemma}
\newtheorem{proposition}[theorem]{Proposition}
\theoremstyle{remark}
\newtheorem{remark}[theorem]{Remark}

\title{The Integer $\{2\}$-Domination Number of Grids\\ with at Most Ten Rows}
\author{Camden Link}
\date{Draft of November 22, 2026}

\begin{document}
\maketitle

\begin{abstract}
An integer $\{2\}$-dominating function of a graph assigns to each vertex a weight $0$, $1$ or
$2$ so that the weights on every vertex and its neighbours sum to at least $2$; the integer
$\{2\}$-domination number is the least total weight. Lee and Liu determined this number for
grids with one and two rows and conjectured formulas for three and four rows. Both conjectures
are true; this was first shown in an online repository of machine-checked proofs, and we give
independent proofs. Our main results are new: for every number of rows from five to ten we
determine the integer $\{2\}$-domination number of the grid for every number of columns. For
nine rows it is exactly $4n+2$. The proofs treat the problem as a shortest-path problem over
column states in the $(\min,+)$ algebra. For up to five rows the results are formally verified
in the Lean proof assistant, by certificates that Lean's kernel evaluates; for six to ten rows
they rest on exact computation, cross-checked as described.
\end{abstract}

\section{Introduction}

Let $G=(V,E)$ be a finite graph and write $N[v]$ for the closed neighbourhood of a vertex $v$.
For a positive integer $k$, an \emph{integer $\{k\}$-dominating function} of $G$ is a function
$f\colon V\to\{0,1,2,\dots\}$ with $\sum_{u\in N[v]}f(u)\ge k$ for every $v\in V$, and the
\emph{integer $\{k\}$-domination number} $\gamma_{\{k\}}(G)$ is the least value of
$\sum_{v\in V}f(v)$ over such functions. The notion is due to Domke, Hedetniemi, Laskar and
Fricke~\cite{dhlf}; see also Bre\v{s}ar, Henning and Klav\v{z}ar~\cite{bhk}. For $k=1$ it is the
ordinary domination number $\gamma(G)$. Replacing every value of $f$ above $k$ by $k$ keeps the
condition, so one may take $f\colon V\to\{0,\dots,k\}$. The number $\g(G)$ should not be
confused with the $2$-domination, double domination, Roman or Italian domination numbers; the
known relations among these parameters are inequalities~\cite{bbgms}.

We study the grid $\G mn$, the Cartesian product of paths with $m$ and $n$ vertices, drawn with
$m$ rows and $n$ columns, and write
\[
\gamma(m,n)=\g(\G mn).
\]
Clearly $\gamma(m,n)=\gamma(n,m)$. Lee and Liu~\cite{ll} proved that $\gamma(1,n)=2\ce{n/3}$ and
$\gamma(2,n)=n+1$, proved that $\gamma(3,n)\le\fl{3n/2}+1$, described an algorithm that computes
$\gamma(m,n)$ column by column, listed $\gamma(4,n)$ for $n\le100$, and made two conjectures.

\begin{theorem}[Conjectures 2.5 and 3.7 of Lee and Liu]\label{thm:ll}
\begin{enumerate}
\item[(a)] $\gamma(3,n)=\fl{3n/2}+1$ for every $n\ge1$.
\item[(b)] $\gamma(4,n)=2n-\fl{(n-9)/10}$ for every $n\ge49$. More precisely, for every $n\ge1$,
\[
\gamma(4,n)=\Bigl\lceil\frac{19n+9}{10}\Bigr\rceil+\varepsilon(n),
\qquad
\varepsilon(n)=\begin{cases}1&\text{if }n=1,\\-1&\text{if }n\in\{4,7,8,15,18,26,37,48\},\\0&\text{otherwise.}\end{cases}
\]
\end{enumerate}
\end{theorem}

For $n\ge9$ the expressions $\ce{(19n+9)/10}$ and $2n-\fl{(n-9)/10}$ coincide, and
$\gamma(4,48)=92$ is one less than the formula, so the bound $49$ in (b) is best possible.

Theorem~\ref{thm:ll} is not new. A proof of (a) formalized in Lean was merged into the public
repository~\cite{omega} on 18 October 2026, and a proof of (b) by exact computation, of the kind
described in Section~\ref{sec:periodic}, was posted there on 16 October 2026; that repository
consists of machine-generated, machine-checked mathematics and its results have not appeared in
a journal. We learned of it from a literature search made while the present proofs were being
formalized, and did not use it. We include Theorem~\ref{thm:ll} because our proofs and
formalization are independent, because the Lean statement for four rows covers every $n\ge1$,
and because the same method gives the following. We found no earlier statement of these values;
however, Hare~\cite{hare90} studied the same parameter for products of paths in 1990, under the
name $k$-weight domination, in a paper that we were not able to consult, so we cannot exclude
that some of them appear there.

\begin{theorem}\label{thm:main}
\begin{enumerate}
\item[(5)] $\gamma(5,n)=\ce{(23n+17)/10}$ for every $n\ge1$, except that it is one less for
$n\in\{2,3,5,6,8,9,12,15,18\}$. Equivalently,
$\gamma(5,n)=\min\bigl\{\fl{(7n+5)/3},\,\ce{(23n+17)/10}\bigr\}$ for every $n\ge1$.
\item[(6)] $\gamma(6,n)=\ce{(41n+18)/15}$ for every $n\ge1$, except that $\gamma(6,16)=46$ is
one more.
\item[(7)] $\gamma(7,n)=3n+3+\fl{n/11}+\delta(n)$ for every $n\ge1$, where $\delta(n)=1$ if
$n\equiv8$ or $10\pmod{11}$ and $\delta(n)=0$ otherwise, except that it is one less for
$n\in\{2,3,4,8\}$.
\item[(8)] $\gamma(8,n)=\ce{(46n+32)/13}$ for every $n\ge1$, except that it is one less for
$n\in\{2,3,4,5,7,9,16,18\}$.
\item[(9)] $\gamma(9,n)=4n+2$ for every $n\ge1$.
\item[(10)] $\gamma(10,n)=\ce{(57n+28)/13}+\eta(n)$ for every $n\ge52$, where $\eta(n)=-1$
if $n\equiv1$, $5$ or $21\pmod{26}$, $\eta(n)=1$ if $n\equiv2\pmod{26}$, and $\eta(n)=0$
otherwise. For $1\le n\le51$ the same formula holds except that the value is larger by $1$ for
$n\in\{5,7,10,17,21,25,31,33,36,51\}$, larger by $2$ for $n=1$, and smaller by $1$ for $n=2$.
\end{enumerate}
\end{theorem}

Table~\ref{tab:summary} summarizes the results and how each was verified. Parts (a), (b) and
(5) are proved in Lean (Section~\ref{sec:lean}). Parts (6) to (9) are proved by exact
computation with two independent programs, and part (10) by one program
(Section~\ref{sec:comp}).

\begin{table}[ht]
\centering
\begin{tabular}{rrcccl}
\toprule
$m$ & states & growth per column & $N$ & $p$ & verification\\
\midrule
1 & 6 & $2/3$ & 3 & 3 & Lee and Liu~\cite{ll}\\
2 & 23 & $1$ & 5 & 2 & Lee and Liu~\cite{ll}; Lean\\
3 & 95 & $3/2$ & 15 & 4 & Lean\\
4 & 389 & $19/10$ & 67 & 10 & Lean\\
5 & 1595 & $23/10$ & 30 & 10 & Lean\\
6 & 6538 & $41/15$ & 30 & 15 & two programs\\
7 & 26802 & $34/11$ & 24 & 11 & two programs\\
8 & 109872 & $46/13$ & 34 & 13 & two programs\\
9 & 450410 & $4$ & 38 & 3 & two programs\\
10 & 1846415 & $57/13$ & 61 & 26 & one program\\
\bottomrule
\end{tabular}
\caption{For each number of rows $m$: the number of states of the column automaton of
Section~\ref{sec:automaton}, the limit of $\gamma(m,n)/n$, and the numbers $N$ and $p$ of
Proposition~\ref{prop:periodic}, for which the vector of the dynamic programme repeats with
period $p$ from column $N$ on.}\label{tab:summary}
\end{table}

For $m\le3$ the growth per column is exactly twice that of the domination number, so
$\gamma(m,n)$ stays within a bounded distance of the trivial bound $2\gamma(\G mn)$. From four
rows on it is smaller. We return to this in Section~\ref{sec:open}.

The method is the transfer-matrix method over the $(\min,+)$ semiring, which goes back to
Klav\v{z}ar and \v{Z}erovnik~\cite{kz} and Livingston and Stout~\cite{ls} for domination-type
problems on grids of fixed height, and was used by Gon\c{c}alves, Pinlou, Rao and
Thomass\'e~\cite{gprt} and by Rao and Talon~\cite{rt} for the domination, $2$-domination and
Roman domination numbers of all grids. Livingston and Stout observed that, for a fixed number
of rows, the dynamic programme is eventually periodic, so that finitely many steps determine a
formula for all $n$. Lee and Liu's algorithm computes the same dynamic programme, but for one
$n$ at a time, which is why their statements for three and four rows remained conjectures.
Section~\ref{sec:cert} puts the argument in a form that a proof assistant can check without
trusting the program that found it.

\section{The column automaton}\label{sec:automaton}

Fix $m$. A \emph{column} is a vector $c=(c_1,\dots,c_m)\in\{0,1,2\}^m$, with weight
$|c|=c_1+\dots+c_m$; put $c_0=c_{m+1}=0$ and
\[
\tau_i(c)=c_{i-1}+c_i+c_{i+1}\qquad(1\le i\le m),
\]
the weight that the vertex in row $i$ of the column receives from the column itself. A function
$f\colon V(\G mn)\to\{0,1,2\}$ is a sequence of columns $c^1,\dots,c^n$; put
$c^0=c^{n+1}=0$. The vertex in row $i$ of column $j$ satisfies the domination condition if and
only if
\begin{equation}\label{eq:local}
c^{j-1}_i+\tau_i(c^j)+c^{j+1}_i\ge2 .
\end{equation}
For columns $a,b$ define the \emph{need} vector $r(a,b)$ by
$r(a,b)_i=\max\bigl(0,\,2-a_i-\tau_i(b)\bigr)$. Then \eqref{eq:local} holds for all rows $i$ of
column $j$ if and only if $c^{j+1}\ge r(c^{j-1},c^j)$ in every coordinate.

Thus, when the columns are read from left to right, all that matters about the past is the pair
$(b,\nu)$ consisting of the last column $b$ and the need vector $\nu$ that the next column must
dominate. Let $S_m$ be the set of pairs $(b,\nu)$ with $0\le\nu_i\le\max(0,2-\tau_i(b))$ for all
$i$; these are exactly the pairs $(b,r(a,b))$. Consider the automaton with state set $S_m$,
initial state $(0,0)$, a transition
\[
(b,\nu)\ \xrightarrow{\;c\;}\ (c,\,r(b,c))\qquad\text{whenever } c\ge\nu\text{ coordinatewise},
\]
of cost $|c|$, and final states the pairs $(b,0)$.

\begin{lemma}\label{lem:bridge}
The integer $\{2\}$-dominating functions $f\colon V(\G mn)\to\{0,1,2\}$ correspond bijectively to
the walks of $n$ transitions from the initial state to a final state, and the weight of $f$ is
the cost of the walk. Hence $\gamma(m,n)$ is the least cost of such a walk.
\end{lemma}

\begin{proof}
Reading $c^1,\dots,c^n$ from the initial state, the transition reading $c^{j+1}$ is allowed if
and only if $c^{j+1}\ge r(c^{j-1},c^j)$, that is, if and only if column $j$ is dominated; this
holds trivially for $j=0$. The last state $(c^n,r(c^{n-1},c^n))$ is final if and only if
$0=c^{n+1}\ge r(c^{n-1},c^n)$, that is, if and only if column $n$ is dominated.
\end{proof}

The numbers $|S_m|$ for $m=1,\dots,10$ are $6$, $23$, $95$, $389$, $1595$, $6538$, $26802$,
$109872$, $450410$, $1846415$. The first seven agree with the numbers of ``feasible columns''
counted by Lee and Liu.

\section{Eventual periodicity}\label{sec:periodic}

For $n\ge0$ and $s\in S_m$ let $v_n(s)$ be the least cost of a walk of $n$ transitions from the
initial state to $s$, or $\infty$ if there is none. Then $v_0$ is $0$ at the initial state and
$\infty$ elsewhere,
\begin{equation}\label{eq:dp}
v_{n+1}(t)=\min\bigl\{\,v_n(s)+|c| \;:\; s\xrightarrow{\,c\,}t\,\bigr\},
\qquad
\gamma(m,n)=\min\{v_n(s): s\text{ final}\}.
\end{equation}
Write \eqref{eq:dp} as $v_{n+1}=T(v_n)$. The map $T$ commutes with adding a constant to every
entry.

\begin{proposition}\label{prop:periodic}
Suppose that for some $N\ge0$, $p\ge1$ and $q$ all entries of $v_N$ are finite and
$v_{N+p}=v_N+q$ entrywise. Then $v_{n+p}=v_n+q$ for every $n\ge N$, and therefore
$\gamma(m,n+p)=\gamma(m,n)+q$ for every $n\ge N$.
\end{proposition}

\begin{proof}
If $v_{n+p}=v_n+q$ then $v_{n+1+p}=T(v_n+q)=T(v_n)+q=v_{n+1}+q$.
\end{proof}

So a single coincidence between two vectors, found by running \eqref{eq:dp}, determines
$\gamma(m,n)$ for all $n$ from its values for $n<N+p$. Such a coincidence must occur: the
transition graph is strongly connected, and the powers of an irreducible matrix over the
$(\min,+)$ semiring are ultimately periodic up to an additive constant~\cite[Chapter~3]{bcoq}.
Table~\ref{tab:summary} lists, for each $m$, values of $N$ and $p$ for which the hypothesis
holds (the first coincidence that the computation meets). With the values of
$\gamma(m,n)$ for $n<N+p$, this proves Theorems~\ref{thm:ll} and~\ref{thm:main} by exact
computation. The repeating block need not have length $p$: for five rows the vector has period
$10$, while every repeating block of columns of optimal density has length a multiple of $20$,
the second half being the mirror image of the first (Figure~\ref{fig:blocks}).

\section{Upper bounds: patterns}\label{sec:upper}

By Lemma~\ref{lem:bridge} an upper bound for $\gamma(m,n)$ is a list of columns. If such a list
contains positions $i<i'$ with $(c^{i},c^{i+1})=(c^{i'},c^{i'+1})$, then the block
$c^{i+1},\dots,c^{i'}$ may be repeated any number of times, because every window of three
consecutive columns in the longer list already occurs in the original one and \eqref{eq:local}
involves only such windows. Figures~\ref{fig:blocks} and~\ref{fig:nine} show repeatable blocks
of optimal density for $3\le m\le 9$. For each $m$, a finite set of lists containing such blocks, one or more for
each residue of $n$ modulo the block length, gives valid lists of every length whose weights
are the values in Theorems~\ref{thm:ll} and~\ref{thm:main}.

\begin{figure}[ht]
\small
\begin{verbatim}
 m=3 (4 columns, weight 6)      m=4 (10 columns, weight 19)
   . 1 . 1                        . . 2 . . 1 . 1 . 1
   . 1 . 1                        1 . . . 1 1 . 1 . 1
   . 1 . 1                        . 1 . 1 . . . 1 . .
                                  . 1 . 1 . 1 1 . 1 1

 m=5 (20 columns, weight 46)
   . 1 . 1 . 1 1 . . 1 1 . 1 . 1 1 . . 1 1
   . 1 . 1 . . . 1 1 . . . 2 . . . 1 1 . .
   1 . . . 1 1 . . . 1 1 . . . 1 1 . . . 1
   . . 2 . . . 1 1 . . . 1 . 1 . . . 1 1 .
   1 . 1 . 1 1 . . 1 1 . 1 . 1 . 1 1 . . 1

 m=6 (15 columns, weight 41)          m=7 (11 columns, weight 34)
   . 1 1 . 1 . 1 . 1 1 . . 1 1 .        . 1 . 1 . 1 1 . . 1 1
   1 . . . 1 . 1 . . . 1 1 . . 1        . 1 . 1 . . . 1 1 . .
   . . 1 1 . . . 1 1 . . . . 2 .        1 . . . 1 1 . . . . 1
   1 1 . . . 2 . . . 1 1 1 . . .        . . 2 . . . 1 1 1 1 .
   . . . 1 1 . . 1 1 . . . 1 . 1        1 . . . 1 1 . . . . 1
   . 1 1 . . 1 1 . . 1 1 . 1 . 1        . 1 . 1 . . . 1 1 . .
                                        . 1 . 1 . 1 1 . . 1 1

 m=8 (26 columns, weight 92)
   . . 1 1 . 1 . 1 . 1 1 . . 1 1 . 1 . 1 . 1 1 . . 1 1
   1 1 . . . 1 . 1 . . . 1 1 . . . 1 . 1 . . . 1 1 . .
   . . . 1 1 . . . 1 1 . . . . 1 1 . . . 1 1 . . . 1 1
   . 1 1 . . . 2 . . . 1 1 1 1 . . . 2 . . . 1 1 . . .
   1 . . . 2 . . . 1 1 . . . . 1 1 . . . 2 . . . 1 1 1
   . 1 1 . . . 1 1 . . . 1 1 . . . 1 1 . . . 1 1 . . .
   . . . 1 . 1 . . . 1 1 . . 1 1 . . . 1 . 1 . . . 1 1
   1 1 . 1 . 1 . 1 1 . . 1 1 . . 1 1 . 1 . 1 . 1 1 . .
\end{verbatim}
\caption{Blocks of columns that can be repeated indefinitely in an optimal pattern. A dot is
weight $0$. The weight divided by the number of columns is the growth rate in
Table~\ref{tab:summary}.}\label{fig:blocks}
\end{figure}

\begin{figure}[ht]
\small
\begin{verbatim}
 m=9 (26 columns, weight 104)
   . . 1 1 . . 2 . . 1 . 1 . 1 . 1 . 1 1 . 1 . 1 . 1 1
   1 1 . . 1 . . . 1 1 . 1 . 1 . 1 . . . . 1 . 1 . . .
   . . 1 . 1 1 . 1 . . . 1 1 . . . 1 1 1 1 . . . 1 1 .
   1 . 1 . . . . 1 . 1 1 . . . 2 . . . . . . 2 . . . 1
   . . . 1 1 1 1 . . . 1 . 1 . . . 1 1 1 1 . . . 1 . 1
   . 2 . . . . . . 2 . . . 1 1 . 1 . . . . 1 . 1 1 . .
   . . . 1 1 1 1 . . . 1 1 . . . 1 . 1 1 . 1 . . . 1 1
   1 . 1 . . . . 1 . 1 . . . 1 1 . . 1 . . . 1 1 . 1 .
   1 . 1 . 1 1 . 1 . 1 . 1 1 . . 1 1 . . 2 . . 1 . 1 .
\end{verbatim}
\caption{A repeatable block for nine rows, of density $4$.}\label{fig:nine}
\end{figure}

Short grids can be cheaper than the long-run formula. For four rows there is a second block,
of $11$ columns and weight $21$. Its density $21/11$ exceeds $19/10$, but it meets the two ends
of the grid at lower cost: it gives patterns of weight $(21n+4)/11$ whenever
$n\equiv4\pmod{11}$, and for $n=4,15,26,37,48$ this is one less than $\ce{(19n+9)/10}$. For
instance
\begin{verbatim}
   . 1 1 . . 1 1 . 1 1 . . 1 1 .
   1 . . 1 1 . . 1 . . 1 1 . . 1
   1 . . 1 . . . 2 . . . 1 . . 1
   . 1 1 . . 2 . . . 2 . . 1 1 .
\end{verbatim}
has weight $29=\gamma(4,15)$. From $n=59$ on the two families cost the same or the first is
cheaper. This is the source of most of the exceptional values in Theorem~\ref{thm:ll}(b) and
of the threshold $49$. Five rows behave in the same way: $\gamma(5,n)=\fl{(7n+5)/3}$ for
$n\le30$, a growth of $7/3$ per column, and the slower growth $23/10$ takes over from $n=31$.

\section{Lower bounds: certificates}\label{sec:cert}

Proposition~\ref{prop:periodic} proves the lower bounds, but checking it means running the
dynamic programme. For the formal proofs we use a statement that can be checked one transition
at a time. Fix positive integers $K,R$ and integers $C$ and $J\ge0$.

\begin{lemma}\label{lem:cert}
Suppose there are functions $L_0,\dots,L_J\colon S_m\to\N\cup\{\infty\}$ and $F\colon S_m\to\Z$
such that
\begin{enumerate}
\item[(i)] $L_0$ is $0$ at the initial state;
\item[(ii)] $L_{j+1}(t)\le L_j(s)+|c|$ for every transition $s\xrightarrow{c}t$ and every $j<J$;
\item[(iii)] $F(s)+RJ\le K\,L_J(s)$ for every state $s$;
\item[(iv)] $F(t)+R\le F(s)+K|c|$ for every transition $s\xrightarrow{c}t$;
\item[(v)] $F(s)\ge C$ for every final state $s$.
\end{enumerate}
Then $K\,\gamma(m,n)\ge Rn+C$ for every $n\ge J$, and $\gamma(m,n)\ge\min\{L_n(s):s\text{ final}\}$
for every $n\le J$.
\end{lemma}

\begin{proof}
By (i) and (ii), $L_j\le v_j$ for $j\le J$, which gives the second statement. By (iii),
$F(s)+RJ\le K v_J(s)$, and by (iv) and induction $F(s)+Rn\le K v_n(s)$ for all $n\ge J$. Apply
(v) to a final state attaining $\gamma(m,n)$.
\end{proof}

Taking $L_j=v_j$ and $F(s)=\inf_{n\ge J}\bigl(Kv_n(s)-Rn\bigr)$ shows that a certificate exists
whenever $K\gamma(m,n)\ge Rn+C$ holds for all $n\ge J$ and $R/K$ is the growth rate. Since
$\gamma(m,n)$ is an integer, the conclusion reads $\gamma(m,n)\ge\ce{(Rn+C)/K}$. The values used
are $(K,R,C,J)=(1,1,1,1)$ for two rows, $(2,3,1,1)$ for three, $(10,19,9,49)$ for four and
$(10,23,17,19)$ for five. For two and three rows the potential $F$ alone gives the exact answer
for all $n\ge1$; for four and five rows the tables $L_0,\dots,L_{49}$ and $L_0,\dots,L_{19}$
are needed because of the exceptional short grids described in Section~\ref{sec:upper}.

\section{Formal verification}\label{sec:lean}

The following are proved in Lean~4~\cite{lean4} with Mathlib~\cite{mathlib}. The grid is
Mathlib's Cartesian product of two path graphs,
\begin{center}
\texttt{SimpleGraph.pathGraph m $\square$ SimpleGraph.pathGraph n},
\end{center}
on the vertex set \texttt{Fin m $\times$ Fin n}. For a graph $G$, \texttt{IsDomFn G k f} says that
$f\colon V\to\N$ satisfies $k\le f(v)+\sum_{u\in N(v)}f(u)$ for every vertex $v$, and
\texttt{domNum G k} is the least value of $\sum_v f(v)$ over such $f$.

\begin{itemize}
\item \texttt{domNum\_grid\_two}: $\gamma(2,n)=n+1$ for $n\ge1$.
\item \texttt{domNum\_grid\_three}: Theorem~\ref{thm:ll}(a).
\item \texttt{domNum\_grid\_four} and \texttt{domNum\_grid\_four\_of\_ge}:
Theorem~\ref{thm:ll}(b), for all $n\ge1$ and in Lee and Liu's form for $n\ge49$;
\texttt{domNum\_grid\_four\_48}: $\gamma(4,48)=92$.
\item \texttt{domNum\_grid\_five} and \texttt{domNum\_grid\_five\_of\_ge}:
Theorem~\ref{thm:main}(5), for all $n\ge1$ and as $\ce{(23n+17)/10}$ for $n\ge19$.
\item \texttt{exists\_dom\_grid\_three}, \texttt{\_four}, \texttt{\_five}: the minimum is
attained by a function with values in $\{0,1,2\}$.
\end{itemize}

The statements are about the graph; the automaton appears only in the proofs. The file
\texttt{Bridge.lean} proves Lemma~\ref{lem:bridge} for all $m$ and $n$, with columns encoded as
numbers in base $3$, and \texttt{Certificate.lean} proves Lemma~\ref{lem:cert} and the
repetition argument of Section~\ref{sec:upper}. These two files are written by hand. A program computes
the data: the tables $L_j$ and $F$, a transition table, and the lists of columns for the upper
bounds. The generated files contain this data and assertions that certain boolean functions
return \texttt{true} on it, proved by evaluation in Lean's kernel. An error in the program can
therefore cause a proof to fail but cannot cause a false statement to be accepted. No
\texttt{native\_decide} is used, and all proofs use no axioms beyond the three standard ones of
Lean's classical logic.

Kernel evaluation is slow compared with compiled code, and two points made the computation
feasible. Tables are packed into single natural numbers, so that a lookup is a shift and a
remainder, which the kernel performs with its built-in arithmetic. And the transitions are not
recomputed in each round of condition~(ii): a transition table is checked once against the
definition of the automaton, and conditions (ii) and (iv) then read it. For five rows the
table has $1595$ rows of $243$ entries, the conditions are checked in $236$ separate kernel
computations, and the generated file has $1.3$ megabytes; checking it takes about ten minutes
on a laptop, and the whole development about twelve. The checks have to be run one at a time:
by default Lean elaborates the proofs in a file in parallel, and the kernel's caches for all of
them together exhaust the memory of a small machine.
Six rows would need a table about fifty times larger and was not attempted.

\section{Computation for six to ten rows}\label{sec:comp}

For $m\ge6$ the results rest on exact integer computation by the following programs, which are
in the repository with the Lean development.

\emph{Lower and upper bounds together.} A program runs \eqref{eq:dp} until it finds $N$ and $p$
with $v_{N+p}=v_N+q$. For each last column $b$ it computes
$\min\{v_n(b,\nu):\nu\le c\}$ for all $c$ at once by running minima along each coordinate, so a
step costs about $m\,9^m$ operations instead of one operation per transition. Ten rows
($1\,846\,415$ states) took about two hours of processor time on a laptop.

\emph{Independent check.} A second program uses a different dynamic programme, whose states are
all $9^m$ pairs of consecutive columns with no compression, and the dual transform (minima over
upper sets). For $3\le m\le9$ it reproduces every value $\gamma(m,n)$ for $n\le N+p$ and
verifies the periodicity $V_{N+1+p}=V_{N+1}+q$ of its own vector. For ten rows it would need
more memory than was available, so for that height the lower bound depends on the first program
alone.

\emph{Upper bounds from the definition.} The first program also prints an optimal list of
columns for each $n$ up to a bound beyond $N+p$. A short script rebuilds each list as an
$m\times n$ array, checks the domination condition at every vertex directly, finds repeatable
blocks as in Section~\ref{sec:upper}, and confirms that the lists and their extensions cover
every $n\ge1$ with the weights of Theorems~\ref{thm:ll} and~\ref{thm:main}. This was done for
$1\le m\le 9$; it does not depend on any dynamic programme. For ten rows the lists for
$n\le125$ were checked in the same way; their repeatable blocks have $52$ columns and do not
cover every residue of $n$, so for larger $n$ the upper bound in part (10) rests on the first
program.

\emph{Further checks.} The identity $\gamma(m,n)=\gamma(n,m)$ holds for all computed values
with $m,n\le10$, where the two sides come from automata of different sizes. Exhaustive search
over all functions agrees with the dynamic programme for the $28$ grids with at most $14$
vertices and $m\le4$. The values for $m=4$ and $n\le100$ agree with the table of Lee and Liu.

\section{Remarks and open questions}\label{sec:open}

\begin{enumerate}
\item \emph{Comparison with domination and fractional domination.} For every graph,
$2\gamma_f(G)\le\g(G)\le2\gamma(G)$, where $\gamma_f$ is the fractional domination
number~\cite{dhlf,bhk}. Table~\ref{tab:compare} compares the growth per column of the three
quantities for grids, using the bounds of Hare~\cite{hare94} for $\gamma_f$ and the known
formulas for $\gamma$~\cite{aciop,gprt}. For three rows $\g$ grows like $2\gamma$; from four
rows on it lies strictly between the two bounds, so neither determines it.
\begin{table}[ht]
\centering
\begin{tabular}{rccc}
\toprule
$m$ & $2\gamma_f$ & $\g$ & $2\gamma$\\
\midrule
3 & $10/7\approx1.429$ & $3/2=1.5$ & $3/2=1.5$\\
4 & $20/11\approx1.818$ & $19/10=1.9$ & $2$\\
5 & $20/9\approx2.222$ & $23/10=2.3$ & $12/5=2.4$\\
6 & $76/29\approx2.621$ & $41/15\approx2.733$ & $20/7\approx2.857$\\
7 & $142/47\approx3.021$ & $34/11\approx3.091$ & $10/3\approx3.333$\\
8 & $65/19\approx3.421$ & $46/13\approx3.538$ & $15/4=3.75$\\
\bottomrule
\end{tabular}
\caption{Growth per column, as $n\to\infty$, of twice the fractional domination number, of the
integer $\{2\}$-domination number, and of twice the domination number of $\G mn$.}\label{tab:compare}
\end{table}
\item \emph{All grids.} Since every closed neighbourhood has at most five vertices,
$\gamma(m,n)\ge 2mn/5$, and $\gamma(m,n)\le2\gamma(\G mn)$, so $\gamma(m,n)/(mn)\to2/5$ as
$m,n\to\infty$. The growth rates in Table~\ref{tab:summary}, divided by $m$, are
$0.667$, $0.5$, $0.5$, $0.475$, $0.46$, $0.456$, $0.442$, $0.442$, $0.444$, $0.438$ for
$m=1,\dots,10$.
Is there a formula for $\gamma(m,n)$ valid for all sufficiently large $m$
and $n$, in the way that $\gamma(\G mn)=\fl{(m+2)(n+2)/5}-4$ for $16\le m\le n$~\cite{gprt}?
The method of~\cite{gprt,rt}, which bounds the loss at the border of the grid, seems the
natural approach.
\item \emph{Nine rows.} Is there a proof of $\gamma(9,n)=4n+2$ that a person can check by hand?
\item \emph{Square grids.} The values of $\gamma(n,n)$ for $n=1,\dots,10$ are
$2,3,5,8,13,18,24,31,38,47$.
\item \emph{Formal proofs for more rows.} The certificate of Lemma~\ref{lem:cert} for six rows
has $6538$ states and $729$ columns. A more economical encoding, for instance one that uses the
symmetry exchanging top and bottom, might bring it within reach of kernel evaluation.
\item \emph{Other values of $k$.} The automaton extends to integer $\{k\}$-domination with
weights in $\{0,\dots,k\}$; the number of states grows quickly with $k$.
\end{enumerate}

\section{Use of AI tools}\label{sec:ai}

This work was carried out with an AI system, Claude (Anthropic), under the direction of the
author. The AI system read the source paper, carried out the literature search, found the
proofs, wrote the Lean formalization and all computer programs, ran the computations, and wrote
the text of this draft. The author chose the problem from a list, directed the work, and
decided what to pursue. The statements listed in Section~\ref{sec:lean} have been checked by
Lean; the remaining results have been checked by exact computation as described in
Section~\ref{sec:comp} but have not been reviewed by a mathematician. The literature search,
including what could not be accessed, is recorded in the repository.

\begin{thebibliography}{99}

\bibitem{aciop}
S. Alanko, S. Crevals, A. Isopoussu, P. \"Osterg{\aa}rd and V. Pettersson, Computing the
domination number of grid graphs, \emph{Electron. J. Combin.} \textbf{18} (2011), Paper~141.

\bibitem{bcoq}
F. Baccelli, G. Cohen, G. J. Olsder and J.-P. Quadrat, \emph{Synchronization and Linearity: An
Algebra for Discrete Event Systems}, Wiley, 1992.

\bibitem{bbgms}
F. Bonomo, B. Bre\v{s}ar, L. N. Grippo, M. Milani\v{c} and M. D. Safe, Domination parameters
with number $2$: interrelations and algorithmic consequences, \emph{Discrete Appl. Math.}
\textbf{235} (2018), 23--50.

\bibitem{bhk}
B. Bre\v{s}ar, M. A. Henning and S. Klav\v{z}ar, On integer domination in graphs and
Vizing-like problems, \emph{Taiwanese J. Math.} \textbf{10} (2006), 1317--1328.

\bibitem{dhlf}
G. S. Domke, S. T. Hedetniemi, R. C. Laskar and G. Fricke, Relationships between integer and
fractional parameters of graphs, in \emph{Graph Theory, Combinatorics, and Applications},
Vol.~1, Wiley, 1991, pp.~371--387.

\bibitem{gprt}
D. Gon\c{c}alves, A. Pinlou, M. Rao and S. Thomass\'e, The domination number of grids,
\emph{SIAM J. Discrete Math.} \textbf{25} (2011), 1443--1453.

\bibitem{hare90}
E. O. Hare, $k$-weight domination and fractional domination of $P_m\times P_n$, \emph{Congr.
Numer.} \textbf{78} (1990), 71--80.

\bibitem{hare94}
E. O. Hare, Fibonacci numbers and fractional domination of $P_m\times P_n$, \emph{Fibonacci
Quart.} \textbf{32} (1994), 69--73.

\bibitem{kz}
S. Klav\v{z}ar and J. \v{Z}erovnik, Algebraic approach to fasciagraphs and rotagraphs,
\emph{Discrete Appl. Math.} \textbf{68} (1996), 93--100.

\bibitem{ll}
J.-Y. Lee and C.-A. Liu, The integer $\{2\}$-domination number of grids, preprint, 2025,
\url{https://arxiv.org/abs/2502.00134}.

\bibitem{ls}
M. Livingston and Q. F. Stout, Constant time computation of minimum dominating sets,
\emph{Congr. Numer.} \textbf{105} (1994), 116--128.

\bibitem{mathlib}
The mathlib Community, The Lean mathematical library, in \emph{Proceedings of the 9th ACM
SIGPLAN International Conference on Certified Programs and Proofs}, ACM, 2020, pp.~367--381.

\bibitem{lean4}
L. de Moura and S. Ullrich, The Lean 4 theorem prover and programming language, in
\emph{Automated Deduction -- CADE 28}, Lecture Notes in Comput. Sci., Vol.~12699, Springer,
2021, pp.~625--635.

\bibitem{omega}
The Omega Institute, \emph{trureturing}, software repository, issues 17651 and 17664 and pull
request 18785, October 2026, \url{https://github.com/the-omega-institute/trureturing}.

\bibitem{rt}
M. Rao and A. Talon, The $2$-domination and Roman domination numbers of grid graphs,
\emph{Discrete Math. Theor. Comput. Sci.} \textbf{21} (2019), no.~1, Paper~9.

\end{thebibliography}

\bigskip
\noindent 2020 \emph{Mathematics Subject Classification}: Primary 05C69; Secondary 05C85, 68V20.

\noindent \emph{Keywords}: integer domination, $\{k\}$-domination, grid graph, $(\min,+)$
algebra, transfer matrix, formal proof.

\end{document}
